\documentclass[11pt,a4paper]{article}
\usepackage[T1]{fontenc}
\usepackage[utf8]{inputenc}
\usepackage{lmodern,amsmath,amssymb}
\usepackage[margin=25mm]{geometry}
\usepackage[hidelinks]{hyperref}
\hypersetup{pdftitle={The one-loop potential along the abelian valley of the torus is periodic, reproduces C133'},pdfauthor={navstokgap research working notes}}
\newcommand{\tightlist}{\setlength{\itemsep}{0pt}\setlength{\parskip}{0pt}}
\setlength{\emergencystretch}{3em}
\setcounter{secnumdepth}{0}
\begin{document}
\hypertarget{the-one-loop-potential-along-the-abelian-valley-of-the-torus-is-periodic-reproduces-c133s-linear-term-at-the-origin-and-saturates-at-order-1l}{%
\section{The one-loop potential along the abelian valley of the torus is
periodic, reproduces C133's linear term at the origin, and saturates at
order
1/L}\label{the-one-loop-potential-along-the-abelian-valley-of-the-torus-is-periodic-reproduces-c133s-linear-term-at-the-origin-and-saturates-at-order-1l}}

For \(SU(2)\) Yang--Mills on \(T^3_L\) with a constant abelian
background \(a=a_i^{\,3}T^3\), the zero-point energy of all modes,
relative to \(a=0\), is the exact periodic function
\[U(a)=\frac{2}{\pi^2L}\,\Phi(La),\qquad
\Phi(b)=\sum_{m\in\mathbb Z^3\setminus\{0\}}\frac{1-\cos(m\cdot b)}{|m|^4}\ \ge0,\]
in units \(\hbar=c=1\), obtained by Poisson summation of the charged
frequencies \(|k\pm a|\) over \(k\in(2\pi/L)\mathbb Z^3\); the divergent
part of the sum is \(a\)-independent under any periodic regularization
and drops out. Near the origin
\(\Phi(b)=\pi^2|b|+\tfrac16C_M|b|^2+O(|b|^3)\) with
\(C_M=\lim_{b\to0}[\sum'_me^{im\cdot b}|m|^{-2}-2\pi^2/|b|]\), so
\(U(a)=2|a|+\tfrac{C_M}{3\pi^2}L|a|^2+\cdots\). The linear term \(2|a|\)
is exactly the \(k=0\) contribution, four oscillators of frequency
\(|a|\) with zero-point energy \(\tfrac12|a|\) each, which are the
transverse oscillators of C133 along its abelian valley: the
field-theory computation reproduces the mechanism of
\href{low-dimensional-mass-gap.md}{G07} and
\href{action-floor-yang-mills-gap.md}{G08} term by term. The nonzero
modes contribute the correction \(U_\perp=U-2|a|=O(L|a|^2)\), of order
\(g^{4/3}/L\) in the rescaled variables of
\href{weak-coupling-feshbach-reduction.md}{the Feshbach note}, which is
the input (H3) needs, with the sign of \(C_M\) (negative for the cubic
lattice) deciding that the correction lowers the valley potential.
Globally \(U\) is bounded by \(2\max\Phi/(\pi^2L)=O(1/L)\) and vanishes
exactly at \(La\in2\pi\mathbb Z^3\), the holonomies that are central in
the fundamental representation: the confinement of the constant modes
along the valley is linear only for \(La\ll1\), saturates at order
\(1/L\), and has degenerate minima at the centre elements, which is the
small-volume origin of the electric-flux sectors. Since the zero-mode
gap is \(g^{2/3}/L\ll1/L\), the ground state is localized at
\(La\sim g^{2/3}\) and never reaches the saturation; the crossover
\(z\simeq2\) is where it does. Sources: this is the spatial-torus form
of Weiss's finite-temperature potential (Phys. Rev.~D 24 (1981) 475,
metadata), and it appears in the small-volume literature (Lüscher 1983,
abstract); the derivation below is self-contained. Nothing here is
promoted.

\hypertarget{frequencies-in-a-constant-abelian-background}{%
\subsection{1. Frequencies in a constant abelian
background}\label{frequencies-in-a-constant-abelian-background}}

Let \(a_i=\alpha_iT^3\), \(i=1,2,3\), with \(\vec\alpha\in\mathbb R^3\),
and write \(a\) for \(\vec\alpha\) when no confusion arises. Decompose
the adjoint fluctuation
\(\tilde A_i=\tilde A_i^3T^3+\tilde A_i^+T^++\tilde A_i^-T^-\). The
covariant derivative \(D_j=\partial_j+[a_j,\cdot]\) acts on the neutral
component as \(\partial_j\) and on the charged components as
\(\partial_j\pm i\alpha_j\), so on the Fourier mode \(e^{ik\cdot x}\),
\(k\in(2\pi/L)\mathbb Z^3\), the charged components have covariant
momentum \(k\pm a\). The quadratic form
\(\frac1{2g^2}\int|D_i\tilde A_j-D_j\tilde A_i|^2\) is, for a constant
abelian background, the free transverse form with \(k\) replaced by
\(k\pm a\) on the charged components, so after the usual transverse
gauge choice the frequencies are
\[\omega^{(0)}_k=|k|\ \ (\text{neutral, two polarizations}),\qquad
\omega^{(\pm)}_k=|k\pm a|\ \ (\text{charged, two polarizations each}),\]
for \(k\ne0\). At \(k=0\) the neutral components are the valley
directions themselves and the charged components of \(\tilde A_j\) have
frequency \(|a|\) for each of the two charges and, since there is no
transversality at \(k=0\), for the components of \(\tilde A_j\)
orthogonal to the valley in field space: with \(a=\alpha_1T^3\) along
\(x_1\) only, these are the charged parts of \(\tilde A_2,\tilde A_3\),
four real oscillators of frequency \(|a|\), while the charged parts of
\(\tilde A_1\) are gauge directions of frequency zero removed by Gauss's
law.

\textbf{The \(k=0\) term is C133.} The SU(2) matrix model of C133 with
\(\vec x_1=\lambda\hat n\), \(\vec x_2=\vec x_3=0\) has potential
\(\tfrac12g^2(|\vec x_1\times\vec x_2|^2+|\vec x_1\times\vec x_3|^2+\cdots)\),
whose transverse Hessian at that point gives the components of
\(\vec x_2,\vec x_3\) orthogonal to \(\hat n\), four oscillators of
frequency \(g|\lambda|/\sqrt m\); in the field-theory variables of
\href{low-dimensional-mass-gap.md}{G07} Proposition 7 this frequency is
\(|a|\). The zero-point energy \(4\cdot\tfrac12|a|=2|a|\) is the
confining term of G07 Theorem 6(1) and the floor of G08 Theorem 2 along
the valley.

\hypertarget{poisson-summation}{%
\subsection{2. Poisson summation}\label{poisson-summation}}

The zero-point energy of the fluctuations relative to \(a=0\), summing
two polarizations and both charges, is
\[U(a)=\frac12\sum_{k\in(2\pi/L)\mathbb Z^3}\sum_{\rm pol=1,2}\big(|k+a|+|k-a|-2|k|\big)
=\sum_{k\in(2\pi/L)\mathbb Z^3}\big(|k+a|+|k-a|-2|k|\big),\] where the
\(k=0\) term is \(2|a|\) and the neutral modes cancel. Each term is
nonnegative by convexity of \(|\cdot|\) and of order \(|a|^2/|k|\), so
the sum diverges in three dimensions; under a periodic regularization
(lattice dispersion \(\omega_{\rm lat}\) with the exact momentum grid)
the divergent part is \(a\)-independent, because
\(\sum_{k\in\text{grid}}\omega_{\rm lat}(k+a)\) is a periodic function
of \(a\) with period \(2\pi/L\) whose mean over a period is the
\(a\)-independent Brillouin-zone integral. The \(a\)-dependent part is
obtained by Poisson summation. For \(f(k)=|k|\) the distributional
Fourier transform is \(\hat f(x)=\int|k|e^{-ik\cdot x}d^3k=-8\pi/|x|^4\)
for \(x\ne0\) (the formula
\(\widehat{|k|^\alpha}=2^{\alpha+3}\pi^{3/2}\frac{\Gamma((3+\alpha)/2)}{\Gamma(-\alpha/2)}|x|^{-3-\alpha}\)
at \(\alpha=1\), with \(\Gamma(-\tfrac12)=-2\sqrt\pi\)), and
\[\sum_{k\in(2\pi/L)\mathbb Z^3}f(k+a)=\Big(\frac L{2\pi}\Big)^3\sum_{m\in\mathbb Z^3}\hat f(mL)\,e^{im\cdot La}\]
with the \(m=0\) term the divergent, \(a\)-independent one. Hence
\[U(a)=\Big(\frac L{2\pi}\Big)^3\sum_{m\ne0}\Big(-\frac{8\pi}{L^4|m|^4}\Big)\big(e^{im\cdot La}+e^{-im\cdot La}-2\big)
=\frac{2}{\pi^2L}\sum_{m\ne0}\frac{1-\cos(m\cdot La)}{|m|^4}=\frac{2}{\pi^2L}\Phi(La).\]
Lattice corrections change \(\hat f\) only at
\(|x|\lesssim a_{\rm lat}\), while the \(m\ne0\) terms live at
\(|x|=|m|L\ge L\), so they are \(O(a_{\rm lat}^2/L^2)\) relative.

Properties of \(\Phi\): it is nonnegative, periodic with period \(2\pi\)
in each component, invariant under the cubic group, bounded by
\(2\sum'|m|^{-4}\), and zero exactly on \(2\pi\mathbb Z^3\).

\hypertarget{expansion-at-the-origin-and-the-sign-of-the-correction}{%
\subsection{3. Expansion at the origin and the sign of the
correction}\label{expansion-at-the-origin-and-the-sign-of-the-correction}}

\(\Phi\) is not smooth at \(0\). Its Laplacian is the periodic Green's
function \(G(b)=\sum'_me^{im\cdot b}|m|^{-2}\), which has the Coulomb
singularity \(2\pi^2/|b|\) at the origin (from
\(\int d^3m\,e^{im\cdot b}|m|^{-2}\)) and a finite remainder:
\(G(b)=2\pi^2/|b|+C_M+O(|b|^2)\), where \(C_M\) is the constant term of
the periodic Coulomb Green's function of the cubic lattice, the
analytically continued Epstein sum \(\sum'_m|m|^{-2}\). Since
\(\Delta(\pi^2|b|)=2\pi^2/|b|\), the function \(\Phi-\pi^2|b|\) has
Laplacian \(C_M+O(|b|^2)\) near \(0\), and the only quadratic form with
cubic symmetry and Laplacian \(C_M\) is \(\tfrac16C_M|b|^2\); so
\[\Phi(b)=\pi^2|b|+\tfrac16C_M|b|^2+O(|b|^3),\qquad
U(a)=2|a|+\frac{C_M}{3\pi^2}\,L|a|^2+O(L^2|a|^3).\] The linear term
\(2|a|\) is the \(k=0\) contribution of Section 1, so the nonzero modes
contribute
\[U_\perp(a)=U(a)-2|a|=\frac{C_M}{3\pi^2}L|a|^2+O(L^2|a|^3),\] with
\(C_M<0\) for the cubic lattice (lattice-sum tables; the sign is not
re-derived here and affects only the direction of a subleading
correction). In the rescaled variables \(a=g^{2/3}\xi/L\) this is
\(\frac{C_M}{3\pi^2}g^{4/3}|\xi|^2/L\), of relative order \(g^{2/3}\) to
the zero-mode Hamiltonian \(g^{2/3}h_3/L\): the size that (H3) of the
Feshbach reduction assumes, now with its coefficient identified.

\hypertarget{saturation-and-the-centre-degenerate-minima}{%
\subsection{4. Saturation and the centre-degenerate
minima}\label{saturation-and-the-centre-degenerate-minima}}

\(U(a)\le2\max\Phi/(\pi^2L)=O(1/L)\) for all \(a\), while the \(k=0\)
term \(2|a|\) alone grows without bound. The linear confinement of the
constant modes along the abelian valley, which in C133 holds for all
\(|a|\), is in the torus theory the small-\(La\) regime of a bounded
periodic function. The minima of \(U\) are at \(La\in2\pi\mathbb Z^3\):
for \(SU(2)\) with the charged modes of charge \(\pm1\), \(La=2\pi e_i\)
corresponds to the fundamental holonomy \(\exp(2\pi iT^3)=-1\), a centre
element, trivial on all adjoint fields but a distinct state of the
theory; the degenerate minima are connected by tunnelling through the
valley barrier of height \(O(1/L)\), which produces the splittings
between the 't Hooft electric-flux sectors of the small-volume spectrum
(van Baal's programme; not derived here).

Consequence for the small-volume expansion: the zero-mode ground state
has width \(La\sim g^{2/3}\) and energy \(g^{2/3}/L\), far below the
valley saturation scale \(1/L\) and the tunnelling barrier, so the C133
picture is exact at leading order and the corrections enter at relative
order \(g^{2/3}\), as in the Feshbach reduction. The regime where the
zero-mode energy reaches \(1/L\), \(g(L)^{2/3}\sim1\), is where the
holonomy delocalizes over the whole compact valley, the flux sectors
become degenerate and the description changes character: the crossover
\(z\simeq2\) of Lüscher--Münster in this language.

\hypertarget{consequence-for-state}{%
\subsection{5. Consequence for STATE}\label{consequence-for-state}}

(H3)'s one-loop input along the abelian valley is now explicit:
\(U_\perp=\frac{C_M}{3\pi^2}L|a|^2+O(L^2|a|^3)\), bounded by \(O(1/L)\)
globally, periodic with centre-degenerate minima. The same Poisson
computation for non-commuting backgrounds requires the Landau-level
spectrum of the covariant Laplacian with constant field strength and the
Nielsen--Olesen mode, which is the next input; on the abelian valley the
reduction's (H3) reduces to a relatively bounded perturbation of \(h_3\)
by \(-|C_M|g^{4/3}|\xi|^2/(3\pi^2L)\) inside the region \(La\ll1\) and
by a bounded potential outside it, which the confining bound of C133
controls.
\end{document}
